\documentclass{article}
\usepackage[utf8]{inputenc}


%Math
\usepackage{multirow, longtable, amsmath, amssymb, amsthm}
\usepackage{indentfirst, graphicx, tabularx, fancyhdr, xcolor, todonotes}
\usepackage[mathscr]{euscript}

%Formatting
\usepackage[margin = 1 in]{geometry}
\usepackage{hyperref, cite}

%Commands
\newcommand{\Q}{{\mathbb{Q}}}
\newcommand{\F}{{\mathbb{F}}}
\newcommand{\Z}{{\mathbb{Z}}}
\newcommand{\N}{{\mathbb{N}}}
\newcommand{\R}{{\mathbb{R}}}
\newcommand{\C}{{\mathbb{C}}}
\newcommand{ \norm }[1]{{ || {#1}||}}
\newcommand{\abs}[1]{{|{#1}|}}
\newcommand{\cl}[1]{{\overline{#1}}}
\newcommand{\pd}{\partial}
\newcommand{\ip}[2]{ \langle {#1},{#2}\rangle }
\newcommand{\ls}{\lesssim}
\newcommand{\gs}{\gtrsim}

%Result Environment
\newcounter{result}[section]
\newenvironment{myresult}[1][]{\refstepcounter{result}\par\medskip
	\noindent \textbf{[\thesection.\theresult] #1} \rmfamily}{\medskip}

%Proof Environment
\newenvironment{myproof}{%
	\par%
	\medskip
	\leftskip=2em\rightskip=2em%
	\noindent\ignorespaces \textit{Proof:} \newline}{%
	\,\, \newline \textit{Q.E.D.}\par\medskip}




\begin{document}


\title{Math 126 Midterm 2 Practice}
\author{UCB Summer 2026}
\date{July 22, 2026}
\maketitle

Remember that the exam is closed book. You will have 50 minutes to complete the Exam. You may use the front and back of each page. There are 3 total problems, so ensure that you have attempted all 3. Please write your full name and SID below. Good luck!

\vspace{1 cm}

\textbf{Name:} 

\vspace{1 cm}

\textbf{SID:}

\vspace{1 cm}

\textbf{Honor Pledge:}

\vspace{1 cm}

\textbf{Formulas:}
Characteristic ODEs: 
\begin{align*}
	x'(s) = D_p F \\
	z'(s) = D_p F \cdot p \\
	p'(s) = -D_{x} F - D_{z}(F) p
\end{align*}

Divergence: $$\int_{\Omega} \nabla \cdot F dx = \int_{\pd \Omega} F \cdot \eta dS$$

D'Alembert: $$u(t,x) = \frac{1}{2} \left[g(x+ct) + g(x-ct) \right] + \frac{1}{2c} \int_{x-ct}^{x+ct} h(\tau) d\tau $$

$$\mathcal{D}_{t,x} = \{ (s,y) \in \R_{+} \times \R \mid x-c(t-s) \leq y \leq x+c(t-s)\}$$

Inhomogeneous Wave: $$u(t,x) = \frac{1}{2c} \int_{\mathcal{D}_{t,x}} f(s,y) dsdy$$

Poisson: $$u(t,x) = \frac{\pd}{\pd t} \left( \frac{t}{2\pi} \int_{D} \frac{g(t-xy)}{\sqrt{1 - |y|^{2}}} dy\right) + \frac{t}{2\pi} \int_{D} \frac{h(t-xy)}{\sqrt{1-|y|^{2}}} dy$$

Kirchoff: $$ u(t,x) = \frac{\pd}{\pd t} \left(\frac{1}{4 \pi t} \int_{\pd B(x,t)} g(y) dS(y) \right) + \frac{1}{4\pi t} \int_{\pd B(x,t)} h(y) dS(y)$$

Helmholtz 1D: $$\phi_{k}(x) = \sin(\frac{k\pi}{l}x), \hspace{0.5 cm} k \in \Z$$
solve $$-\phi(x) = \lambda \phi \hspace{0.5 cm} \phi(0) = \phi(l) = 0$$ for $\lambda = (\frac{k\pi}{l})
^{2}$.

Helmholtz Solutions on the Torus: $$\phi_{k}(x) = e^{ikx} \hspace{0.5 cm} k \in \Z$$

Heat Separation of Variables on the Torus: $u(t,x) = e^{-k^{2}t}e^{ikx}$ for $k \in \Z$. 


Heat Kernel:$$ H_{t}(x) = (4\pi t)^{-n/2} e^{-\frac{|x|^{2}}{4t}} $$

Heat Solution on $\R^{n}$ with initial condition $f \in C^{0}(\R^{n})$: $$u(t,x) = H_{t}(x) * f(x)$$

Fourier Coefficient: $$c_{k}[f] = \frac{1}{2\pi}\int_{0}^{2\pi} f(x) e^{-ikx} dx$$

Poisson Kernel: 

$$P_{r}(\theta) = \sum_{\Z} r^{|k|}e^{ik\theta}$$

Laplace Solution with boundary condition $g$:
$$u(r,\theta) = P_{r}(\theta)*g(\theta)$$




\newpage


\textbf{Problem 1:} \\

True/False

Part A.) $C_{c}^{\infty}(\Omega)$ is a complete subspace of $L^{1}(\Omega)$ (in the $L^{1}$ norm). 

Part B.) Under the mixed boundary conditions $u(0) = 0$, $u'(l) = 0$ on $(0,l)$, the Laplacian is self-adjoint. 

Part C.) The partial Fourier sums of a function $f \in C^{1}(\mathbb{T})$ converge pointwise to $f$. 

Part D.) If $\sum |c_{k}[f]^{2} k^{4}| < \infty$, then $f$ must be $C^{2}$.  

Part E.) The heat equation may be used to describe the motion of particles. 




\textbf{Problem 2:} \\

Let $u(t,x)$ be the solution to the heat equation on $(0,\infty) \times \R$ given by 
\begin{align*}
	u(t,x) = H_{t}(x) * g(x)
\end{align*}
for $g \in C_{c}^{\infty}(\R)$ so $\text{supp}(g) \subset [-1,1]$ and $ 0 \leq g(x) \leq 1$. Show that 
\begin{align*}
	\lim_{t \to \infty} u(t,x) = 0
\end{align*}
for all $x$.  

\textbf{Extension:} what bounds can we put on $u(t,x)$ compared to $g$?


\textbf{Problem 3:}


Suppose $u,\phi \in C^{2}(\Omega) \cap C^{0}(\cl{\Omega})$ for a bounded domain $\Omega \subset \R^{n}$. Let $u$ be subharmonic and $\phi$ harmonic, so $u|_{\pd \Omega} \leq \phi|_{\pd \Omega}$. Show $u \leq \phi$ on $\Omega$. 



\textbf{Problem 4:} 

Let $g(\theta) = (\theta-\pi)^{2}$ on $[0,2\pi)$. Does the Laplace equation admit a solution on $B(0,1) \subset \R^{2}$ with this boundary data? If so, what is it (your solution will be a sum). 

\textbf{Problem 5:} 

Let $u(t,x)$ be a heat solution on the torus with initial condition $f$, so 
\begin{align*}
	u(t,x) = \sum_{\Z} c_{k}[f] e^{-k^{2}t}e^{ikx}
\end{align*}

Under what condition on $f$ does $u(t,x)$ decay in time? 

\textbf{Problem 6:}

Recall that $\{\frac{1}{\sqrt{2\pi}}e^{ikx}\}$ is a basis for $L^{2}(-\pi,\pi))$. Show that the functions $\phi_{k}(x) = \sqrt{\frac{2}{\pi}} \sin(kx)$ form a basis for $L^{2}((0,\pi))$, assuming the $\phi_{k}$ already have magnitude 1. \textit{Hint:} Any function on $L^{2}((0,\pi))$ may be extended to an odd function on $(-\pi,\pi)$. 




\end{document}
